\documentclass[a4paper,12pt]{article}
\usepackage[T2A]{fontenc}
\usepackage[utf8]{inputenc}
\usepackage[russian]{babel}
\usepackage{amsmath,amssymb,mathtools}
\renewcommand{\sfdefault}{cmss}
\renewcommand{\familydefault}{\sfdefault}
\usepackage{icomma}
\usepackage{geometry}
\geometry{left=23mm,right=23mm,top=22mm,bottom=24mm}
\usepackage{microtype}
\usepackage{tikz}
\usetikzlibrary{arrows.meta,positioning,calc}
\usepackage{tabularx,booktabs}
\usepackage{enumitem}
\usepackage{parskip}
\usepackage{fancyhdr}
\usepackage{setspace}
\usepackage{xcolor}
\usepackage[hidelinks]{hyperref}
\definecolor{accent}{HTML}{008080}
\definecolor{darkaccent}{HTML}{005A5A}
\definecolor{softgray}{HTML}{5A6267}
\definecolor{textgray}{HTML}{22292C}
\color{textgray}
\setstretch{1.20}
\setlength{\parindent}{0pt}
\setlength{\parskip}{7pt}
\setlist[enumerate]{leftmargin=*,itemsep=0.45em,topsep=0.3em}
\setlist[itemize]{leftmargin=*,itemsep=0.28em,topsep=0.3em}
\newcommand{\topic}{Силы при движении бруска под действием наклонной тяги}
\newcommand{\ruleline}{\par\vspace{2pt}\noindent\textcolor{accent}{\rule{\linewidth}{0.55pt}}\par\vspace{1pt}}
\newcommand{\secname}[1]{\section{#1}\ruleline}
\newcommand{\subname}[1]{\vspace{0.3em}\textcolor{darkaccent}{\bfseries #1}\par\vspace{-0.2em}}
\renewcommand{\thesection}{\arabic{section}}
\pagestyle{fancy}
\fancyhf{}
\fancyhead[L]{\small\color{softgray}Брусок и наклонная сила тяги}
\fancyhead[R]{\small\color{softgray}Володя Хачатурян}
\fancyfoot[C]{\small\thepage}
\renewcommand{\headrulewidth}{0.55pt}
\renewcommand{\headrule}{\hbox to\headwidth{\color{accent}\leaders\hrule height \headrulewidth\hfill}}
\setlength{\headheight}{25pt}
\begin{document}
\pagenumbering{gobble}
\thispagestyle{empty}
\begin{center}
\vspace*{20mm}
{\Large\color{darkaccent} ФИЗИКА}\par
\vspace{15mm}
{\LARGE Чек-лист}\par
\vspace{16mm}
{\Large\color{accent}\bfseries Силы и проекции}\par
{\Large\color{accent}\bfseries при наклонной тяге}\par
\vspace{16mm}
{\normalsize Второй закон Ньютона. Трение. Геометрия угла $30^\circ$}\par
\vfill
{\large Володя Хачатурян}\par
\vspace{9mm}
{4 ноября 2026 г.}\par
\vspace{14mm}
{\small Лёвин Артём Александрович}\par
{\small эксклюзивно для Володи Хачатуряна}\par
\vspace{14mm}
\begin{tikzpicture}[x=1cm,y=1cm]
 \draw[accent,line width=0.85pt] (0,0) .. controls (1.05,0.46) and (2.05,-0.45) .. (3.05,0)
 .. controls (4.0,0.44) and (5.15,-0.42) .. (6.2,0)
 .. controls (7.2,0.4) and (8.5,-0.4) .. (9.2,0);
\end{tikzpicture}
\end{center}
\clearpage
\pagenumbering{arabic}
\setcounter{page}{1}

\secname{С чего начинается задача}

\textcolor{darkaccent}{Главная идея.} Когда брусок тянут под углом к горизонту, сила тяги одновременно влияет на движение вдоль поверхности и на силу реакции опоры. Поэтому нужно работать с \textbf{проекциями сил}.

\begin{center}
\begin{tikzpicture}[scale=0.97,>=Stealth,thick]
 \draw[softgray,line width=0.8pt] (-1.9,0)--(5.5,0);
 \foreach \x in {-1.7,-1.3,...,5.3} \draw[softgray!50,line width=0.25pt] (\x,0)--++(-0.16,-0.18);
 \fill[softgray!12] (0,0) rectangle (1.6,0.85);
 \draw[softgray] (0,0) rectangle (1.6,0.85);
 \node at (0.8,0.43) {$m$};
 \coordinate (A) at (1.6,0.85);
 \coordinate (B) at (3.55,0.85);
 \coordinate (C) at (3.55,1.976);
 \draw[->,line width=1.1pt,accent] (A)--(C) node[pos=0.75,above left] {$\vec F$};
 \draw[->,softgray,densely dashed] (A)--(B) node[pos=0.52,below] {$F_x$};
 \draw[->,softgray,densely dashed] (B)--(C) node[midway,right] {$F_y$};
 \draw[accent] ($(A)+(0.55,0)$) arc[start angle=0,end angle=30,radius=0.55];
 \node[accent] at (2.47,1.10) {$30^\circ$};
 \draw[->,line width=1pt] (0.50,0.85)--(0.50,2.08) node[above] {$\vec N$};
 \draw[->,line width=1pt] (0.84,0.36)--(0.84,-1.16) node[below] {$m\vec g$};
 \draw[->,line width=1pt] (0,0.43)--(-1.30,0.43) node[left] {$\vec F_{\text{тр}}$};
 \draw[->,softgray] (4.45,-0.65)--(5.25,-0.65) node[right] {$x$};
 \draw[->,softgray] (4.45,-0.65)--(4.45,0.10) node[above] {$y$};
\end{tikzpicture}
\end{center}

На рисунке брусок движется \textbf{вправо}. Сила трения скольжения направлена влево, сила тяжести~--- вертикально вниз, реакция опоры~--- вверх, перпендикулярно поверхности. Наклонная сила $F$ имеет две проекции: $F_x$ и $F_y$.

\subname{Что проверить перед расчётом}
\begin{itemize}
\item Составлены «Дано» и схема всех сил.
\item Ось $x$ направлена вдоль движения, ось $y$~--- вертикально.
\item Понятно, есть ли ускорение по каждой оси.
\item Отмечено, что угол $30^\circ$ отсчитывается \textbf{от горизонтали}.
\end{itemize}

\secname{Проекции силы при угле 30 градусов}

Сила $F$~--- гипотенуза прямоугольного треугольника, $F_x$ и $F_y$~--- его катеты. Катет, лежащий \textbf{против угла $30^\circ$}, равен половине гипотенузы:
\[
\boxed{F_y=\frac F2.}
\]
По теореме Пифагора:
\[
F^2=F_x^2+F_y^2,
\qquad F_x=\sqrt{F^2-\left(\frac F2\right)^2}
=\boxed{\frac{F\sqrt3}{2}}.
\]

\subname{Пример: сила тяги $F=20$ Н}
\[
F_y=\frac{20}{2}=10\ \text{Н},\qquad
F_x=\sqrt{20^2-10^2}=\sqrt{300}=10\sqrt3\ \text{Н}.
\]

\textcolor{darkaccent}{Памятка.} Если неизвестен второй катет, вспомните теорему Пифагора. Пока достаточно свойства угла $30^\circ$; синусы и косинусы для этих задач не требуются.

\secname{Баланс сил по осям}

По второму закону Ньютона алгебраическая сумма проекций сил на выбранную ось равна произведению массы на проекцию ускорения.

\subname{По горизонтали ($Ox$)}
\[
\boxed{F_x-F_{\text{тр}}=ma.}\tag{1}
\]
Если брусок движется равномерно, $a=0$, поэтому $F_x=F_{\text{тр}}$. Если ускорение направлено вправо, $a>0$; если влево~--- $a<0$.

\subname{По вертикали ($Oy$)}
\[
N+F_y-mg=0
\quad\Longrightarrow\quad
\boxed{N=mg-\frac F2}.\tag{2}
\]
Справа стоит нуль, потому что \textbf{вертикальное ускорение равно нулю}. При этом вертикальные силы продолжают действовать.

\subname{Сила трения скольжения}
\[
\boxed{F_{\text{тр}}=\mu N.}\tag{3}
\]
Здесь $\mu$~--- коэффициент трения скольжения, $N$~--- модуль реакции опоры. Объединяя (2) и (3), получаем:
\[
\boxed{F_{\text{тр}}=\mu\left(mg-\frac F2\right)}.\tag{4}
\]

В задачах занятия принимаем $g=10\ \text{м/с}^2$. Формулы (2)--(4) применяем, пока брусок сохраняет контакт с опорой ($N>0$).

\secname{Три типовых задачи}

\subname{Пример 1. Найти силу трения}
Брусок массой $1$ кг тянут силой $10$ Н под углом $30^\circ$. Коэффициент трения $0{,}4$. Найдите силу трения скольжения.
\[
N=mg-\frac F2=1\cdot10-\frac{10}{2}=5\ \text{Н},
\]
\[
F_{\text{тр}}=\mu N=0{,}4\cdot5=\boxed{2\ \text{Н}}.
\]
\textbf{Ключевой приём:} сначала выразить $N$ из вертикального баланса, затем подставить в формулу трения.

\subname{Пример 2. Найти массу}
Силой $12$ Н под углом $30^\circ$ тянут брусок. $\mu=0{,}2$, сила трения равна $2{,}8$ Н. Найдите массу.
\[
N=\frac{F_{\text{тр}}}{\mu}=\frac{2{,}8}{0{,}2}=14\ \text{Н},
\]
\[
mg=N+F_y=14+\frac{12}{2}=20\ \text{Н},\qquad
m=\frac{20}{10}=\boxed{2\ \text{кг}}.
\]
\textbf{Ключевой приём:} использовать уравнение трения для поиска $N$, затем~--- вертикальный баланс для поиска $m$.

\subname{Пример 3. Найти ускорение}
Брусок массой $5$ кг тянут силой $20$ Н под углом $30^\circ$; $\mu=0{,}2$. Брусок скользит вправо. Найдите проекцию ускорения на ось $x$.
\[
N=5\cdot10-\frac{20}{2}=40\ \text{Н},\qquad
F_{\text{тр}}=0{,}2\cdot40=8\ \text{Н}.
\]
\[
F_x=\sqrt{20^2-10^2}=10\sqrt3\ \text{Н},
\]
\[
a=\frac{F_x-F_{\text{тр}}}{m}
=\frac{10\sqrt3-8}{5}
\approx\boxed{1{,}86\ \text{м/с}^2}.
\]
\textbf{Ключевой приём:} недостающую горизонтальную проекцию найти по теореме Пифагора; затем применить уравнение (1).

\clearpage
\thispagestyle{fancy}
\begin{center}
{\LARGE\color{darkaccent} Алгоритм решения}\par
\vspace{0.7em}
{\small Для бруска, который скользит по горизонтали под действием тяги под углом $30^\circ$}\par
\vspace{1.5em}
\begin{tikzpicture}[
 node distance=11mm,
 >=Stealth,
 flowstep/.style={draw=accent!70,rounded corners=2pt,align=center,text width=11.5cm,minimum height=13mm,inner sep=5pt,font=\normalsize},
 terminus/.style={draw=darkaccent,rounded corners=7pt,align=center,text width=8.4cm,minimum height=11mm,inner sep=4pt},
 arr/.style={->,line width=0.9pt,accent}
]
\node[terminus] (begin) {\textbf{Начало: прочитать условие и записать «Дано»}};
\node[flowstep,below=of begin] (forces) {Нарисовать брусок, все действующие силы,\ оси $x$ и $y$, отметить $30^\circ$};
\node[flowstep,below=of forces] (proj) {Найти проекции наклонной тяги:\ $F_y=F/2$, \quad $F_x=\sqrt{F^2-F_y^2}$};
\node[flowstep,below=of proj] (balance) {Составить систему:\ $F_x-F_{\text{тр}}=ma$, \quad $N+F_y-mg=0$, \quad $F_{\text{тр}}=\mu N$};
\node[flowstep,below=of balance] (solve) {Выразить недостающие величины,\ выполнить подстановки и вычисления};
\node[flowstep,below=of solve] (verify) {Проверить знаки, единицы измерения,\ контакт с опорой ($N>0$) и смысл результата};
\node[terminus,below=of verify] (end) {\textbf{Записать ответ}};
\foreach \x/\y in {begin/forces,forces/proj,proj/balance,balance/solve,solve/verify,verify/end} \draw[arr] (\x.south)--(\y.north);
\end{tikzpicture}
\end{center}
\clearpage

\secname{Типичные ошибки и самопроверка}
\begin{enumerate}
\item \textbf{В уравнение по $x$ подставили всю силу $F$.} Вдоль горизонтали действует её проекция $F_x$.
\item \textbf{Записали $N=mg$.} При тяге вверх под углом вертикальная составляющая $F_y$ частично разгружает опору: $N=mg-F_y$.
\item \textbf{По оси $y$ указали «силы не действуют».} Силы есть, но их проекции в сумме равны нулю, так как $a_y=0$.
\item \textbf{Для ускорения использовали сумму модулей сил.} Проекции имеют знаки: $F_x-F_{\text{тр}}=ma$.
\item \textbf{Забыли, какой катет равен половине гипотенузы.} При угле $30^\circ$ к горизонтали это вертикальная проекция $F_y$.
\item \textbf{Остановились, потому что неизвестно $F_x$.} Используйте теорему Пифагора: $F_x^2=F^2-F_y^2$.
\item \textbf{Получили отрицательное ускорение и заменили его на положительное.} Минус у проекции показывает направление ускорения против положительной оси $x$.
\end{enumerate}

\subname{Контрольный список}
\begin{itemize}
\item[$\square$] Все силы на чертеже подписаны и направлены верно.
\item[$\square$] Проекции $F_x$ и $F_y$ найдены без путаницы.
\item[$\square$] Горизонтальное уравнение содержит $ma$.
\item[$\square$] Вертикальное уравнение равно нулю из-за $a_y=0$.
\item[$\square$] Использовано $F_{\text{тр}}=\mu N$.
\item[$\square$] Все подстановки пронумерованы или понятны по записи.
\item[$\square$] Ответ имеет нужные единицы и физический смысл.
\end{itemize}

\clearpage
\section*{Тренировочный блок · Путь А}
\ruleline
Решайте задания последовательно. В каждом: \textbf{«Дано» $\to$ чертёж $\to$ проекции $\to$ баланс сил $\to$ расчёт $\to$ ответ}. Используйте $g=10\ \text{м/с}^2$. Если сказано, что брусок скользит вправо, трение направлено влево. Все силы тяги образуют угол $30^\circ$ с горизонталью и направлены вверх-вправо.

\subname{Базовый уровень}
\begin{enumerate}
\item На брусок действует сила тяги $F=10$ Н. Найдите её горизонтальную и вертикальную проекции $F_x$, $F_y$.
\item Масса бруска $3$ кг, сила тяги $8$ Н. Найдите силу реакции опоры $N$.
\item Брусок массой $2$ кг скользит вправо. Сила тяги $10$ Н, коэффициент трения $\mu=0{,}3$. Найдите силу трения.
\end{enumerate}

\subname{Средний уровень}
\begin{enumerate}[resume]
\item Брусок массой $4$ кг скользит вправо. Сила тяги $12$ Н, $\mu=0{,}25$. Найдите $N$ и $F_{\text{тр}}$.
\item Сила тяги $16$ Н, коэффициент трения $0{,}2$, сила трения $5{,}6$ Н. Найдите массу бруска.
\item Масса бруска $5$ кг, $F=20$ Н, $\mu=0{,}2$. Брусок скользит вправо. Найдите ускорение.
\item Масса бруска $3$ кг, $F=12$ Н, $\mu=0{,}25$. Брусок скользит вправо. Найдите ускорение.
\end{enumerate}

\subname{Повышенный уровень}
\begin{enumerate}[resume]
\item Масса бруска $4$ кг, сила тяги $14$ Н. При скольжении вправо сила трения равна $6$ Н. Найдите коэффициент трения.
\item Брусок массой $5$ кг скользит вправо с ускорением $2\ \text{м/с}^2$. Коэффициент трения $\mu=0{,}2$. Найдите модуль наклонной силы тяги $F$.
\item Брусок массой $4$ кг скользит вправо под действием силы тяги $12$ Н. Коэффициент трения $\mu=0{,}5$. Найдите проекцию ускорения на ось $x$ и объясните, что означает её знак.
\end{enumerate}

\clearpage
\section*{Ответы к тренировочному блоку}
\ruleline
\renewcommand{\arraystretch}{1.55}
\begin{tabularx}{\linewidth}{@{}c X@{}}
\toprule
\textbf{№} & \textbf{Ответ}\\
\midrule
1 & $F_x=5\sqrt3$ Н, $F_y=5$ Н.\\
2 & $N=26$ Н.\\
3 & $F_{\text{тр}}=4{,}5$ Н.\\
4 & $N=34$ Н, $F_{\text{тр}}=8{,}5$ Н.\\
5 & $m=3{,}6$ кг.\\
6 & $a=\dfrac{10\sqrt3-8}{5}\approx1{,}86\ \text{м/с}^2$.\\
7 & $a=2\sqrt3-2\approx1{,}46\ \text{м/с}^2$.\\
8 & $\mu=\dfrac{2}{11}\approx0{,}18$.\\
9 & $F=\dfrac{200}{5\sqrt3+1}\approx20{,}70$ Н.\\
10 & $a=\dfrac{6\sqrt3-17}{4}\approx-1{,}65\ \text{м/с}^2$; ускорение направлено влево, брусок замедляется при движении вправо.\\
\bottomrule
\end{tabularx}

\vspace{1.5em}
\textcolor{darkaccent}{Итоговая памятка}

\begin{tabularx}{\linewidth}{@{}p{0.32\linewidth} X@{}}
\toprule
\textbf{Сигнал в задаче} & \textbf{Полезный шаг}\\
\midrule
Сила под $30^\circ$ & $F_y=F/2$, затем $F_x$ по теореме Пифагора.\\
Найти трение & $N=mg-F_y$, затем $F_{\text{тр}}=\mu N$.\\
Найти массу & Используйте $N=F_{\text{тр}}/\mu$ и вертикальный баланс.\\
Найти ускорение & Из $F_x-F_{\text{тр}}=ma$ выразите $a$.\\
Неизвестна сила тяги & Подставьте проекции и формулу трения в баланс по $x$.\\
\bottomrule
\end{tabularx}
\vfill
\begin{center}\small\color{softgray}
Лёвин Артём Александрович · эксклюзивно для Володи Хачатуряна
\end{center}
\end{document}
